\documentclass[11pt]{article}

% ── Page geometry ──────────────────────────────────────────────────
\usepackage[margin=1in]{geometry}

% ── Fonts & encoding ──────────────────────────────────────────────
\usepackage[T1]{fontenc}
\usepackage[utf8]{inputenc}
\usepackage{lmodern}
\usepackage{microtype}

% ── Mathematics ───────────────────────────────────────────────────
\usepackage{amsmath,amssymb,amsthm}
\usepackage{mathtools}
\usepackage{stmaryrd}

% ── Figures, tables, graphics ─────────────────────────────────────
\usepackage{graphicx}
\usepackage{booktabs}
\usepackage{tabularx}
\usepackage{multirow}
\usepackage{float}

% ── Code listings ─────────────────────────────────────────────────
\usepackage{listings}
\usepackage{xcolor}

\definecolor{japlkw}{RGB}{0,100,170}
\definecolor{japlstr}{RGB}{170,55,0}
\definecolor{japlcmt}{RGB}{100,120,100}
\definecolor{japlbg}{RGB}{248,248,248}

\lstdefinelanguage{japl}{
  morekeywords={fn,let,match,with,type,module,import,if,then,else,
    use,own,ref,opaque,trait,impl,deriving,test,property,forall,
    assert,where,signature,foreign,unsafe,loop,while,do,continue,
    packed,bench},
  morekeywords=[2]{Int,Float,Float32,Bool,Char,String,Bytes,Unit,Never,
    Option,Some,None,Result,Ok,Err,List,Map,Set,Pid,True,False,
    Pure,Io,Async,Net,State,Process,Fail,Reply},
  sensitive=true,
  morecomment=[l]{--},
  morestring=[b]",
  literate={->}{{$\rightarrow$}}2 {=>}{{$\Rightarrow$}}2
           {|>}{{$\triangleright$}}2 {>>}{{$\gg$}}2
           {++}{{$\mathbin{+\!\!+}$}}2 {<>}{{$\diamond$}}2,
}

\lstset{
  language=japl,
  basicstyle=\small\ttfamily,
  keywordstyle=\color{japlkw}\bfseries,
  keywordstyle=[2]\color{japlkw!70!black},
  commentstyle=\color{japlcmt}\itshape,
  stringstyle=\color{japlstr},
  backgroundcolor=\color{japlbg},
  frame=single,
  framerule=0.4pt,
  rulecolor=\color{black!20},
  numbers=left,
  numberstyle=\tiny\color{gray},
  numbersep=8pt,
  breaklines=true,
  showstringspaces=false,
  tabsize=2,
  xleftmargin=1.5em,
  framexleftmargin=1.5em,
  aboveskip=1em,
  belowskip=1em,
  captionpos=b,
}

% ── Hyperlinks ────────────────────────────────────────────────────
\usepackage[colorlinks=true,linkcolor=blue!60!black,
  citecolor=blue!60!black,urlcolor=blue!60!black]{hyperref}

% ── Bibliography ──────────────────────────────────────────────────
\usepackage[numbers,sort&compress]{natbib}

% ── Theorem environments ─────────────────────────────────────────
\newtheorem{theorem}{Theorem}[section]
\newtheorem{lemma}[theorem]{Lemma}
\newtheorem{proposition}[theorem]{Proposition}
\newtheorem{corollary}[theorem]{Corollary}
\theoremstyle{definition}
\newtheorem{definition}[theorem]{Definition}
\newtheorem{example}[theorem]{Example}
\theoremstyle{remark}
\newtheorem{remark}[theorem]{Remark}

% ── Macros ────────────────────────────────────────────────────────
\newcommand{\Set}{\mathbf{Set}}
\newcommand{\Type}{\mathbf{Type}}
\newcommand{\CCC}{\mathbf{CCC}}
\newcommand{\Hom}{\mathrm{Hom}}
\newcommand{\id}{\mathrm{id}}
\newcommand{\sem}[1]{\llbracket #1 \rrbracket}
\newcommand{\japl}{\textsc{Japl}}
\newcommand{\keyword}[1]{\texttt{\textbf{#1}}}

% ══════════════════════════════════════════════════════════════════
% FRONT MATTER
% ══════════════════════════════════════════════════════════════════

\title{\textbf{Values Are Primary:}\\
  Immutability as a Foundation for\\
  Type-Safe Concurrent Programming in \japl}

\author{
  Matthew Long \\
  \textit{The JAPL Research Collaboration} \\
  \textit{YonedaAI Research Collective} \\
  Chicago, IL \\
  \texttt{matthew@yonedaai.com}
}

\date{March 2026}

\begin{document}

\maketitle

% ══════════════════════════════════════════════════════════════════
\begin{abstract}
% ══════════════════════════════════════════════════════════════════

The dominant paradigm in mainstream programming languages treats
\emph{objects with identity} as the primary organising concept:
mutable cells, reference semantics, and pointer aliasing are the
default.  This paper argues that an alternative foundation---one in
which \emph{values} are primary---yields languages that are simpler to
reason about, safer by construction, and better suited to concurrent
and distributed systems.

We develop this thesis in the context of \japl{}, a \emph{strict},
statically typed, effect-aware functional language that makes immutable
algebraic data types the default representation for all data.  Strict
evaluation---in contrast to Haskell's lazy semantics---ensures that
space usage is predictable and that values are fully evaluated, which
simplifies both reasoning and garbage collection.  Mutation is available
but confined to an explicitly tracked resource layer.  We present a
formal framework grounded in category theory---modelling types as
objects of a Cartesian closed category and algebraic data types as
initial algebras of polynomial endofunctors---and give a denotational
semantics for a value-centric core calculus that guarantees referential
transparency.  We describe \japl{}'s type system (parametric
polymorphism, bidirectional inference, row polymorphism for extensible
records, opaque types) and show how the value/resource split
enables a dual memory management strategy: a simplified garbage
collector for immutable data and linear-type--governed ownership for
mutable resources.

Through detailed comparison with Haskell, Rust, Erlang, OCaml, Go,
Gleam, Clojure, and Elixir, we demonstrate that \japl{}'s value model
occupies a distinctive point in the design space.  We evaluate
implementation strategies---hash-array mapped tries, structural
sharing, copy-on-write, packed tagged unions---and present benchmarks
showing that value semantics need not sacrifice performance relative to
mutable approaches.  Finally, we discuss the profound implications of
default immutability for concurrency: values can be freely shared
across process boundaries without synchronization, making \japl{}'s
Erlang-style process model both safe and efficient.

\medskip
\noindent\textbf{Keywords:} value semantics, immutability, algebraic
data types, persistent data structures, type theory, category theory,
concurrency, functional programming, language design.

\end{abstract}

\newpage
\tableofcontents
\newpage

% ══════════════════════════════════════════════════════════════════
\section{Introduction}
\label{sec:intro}
% ══════════════════════════════════════════════════════════════════

\subsection{The Identity Crisis in Programming}

The history of programming languages can be read as a long negotiation
between two worldviews.  In the \emph{identity-centric} view---
exemplified by Simula~\cite{dahl1966simula}, Smalltalk~\cite{kay1993early},
Java~\cite{gosling2000java}, and their descendants---the fundamental
concept is the \emph{object}: an entity with mutable state, behaviour,
and identity.  Two objects may have identical field values yet be
distinct because they occupy different locations in memory.  Programs
are understood as networks of collaborating objects that communicate by
sending messages (method calls) that mutate their internal state.

In the \emph{value-centric} view---rooted in the lambda
calculus~\cite{church1936unsolvable}, refined in ML~\cite{milner1978theory},
Haskell~\cite{jones2003haskell}, and Erlang~\cite{armstrong2007programming}---the
fundamental concept is the \emph{value}: a datum that simply \emph{is}
what it is, with no notion of mutable state or location-dependent
identity.  The number $42$ is $42$ regardless of where it is stored;
the string \texttt{"hello"} does not change over time.  Programs are
understood as compositions of functions that transform values into
other values.

\subsection{The Cost of Identity}

The identity-centric model imposes substantial costs that become
especially acute in concurrent and distributed settings:

\begin{enumerate}
  \item \textbf{Aliasing and mutation.}  When multiple references
    point to the same mutable object, a mutation through one alias is
    visible through all others.  This makes local reasoning
    impossible without whole-program alias analysis, which is
    undecidable in the general case~\cite{landi1992undecidability}.

  \item \textbf{Synchronization overhead.}  In concurrent programs,
    shared mutable state must be protected by locks, atomic
    operations, or transactional mechanisms.  Every lock is a
    potential source of deadlock; every unlocked access is a potential
    data race.  The fundamental difficulty is that identity creates
    \emph{contention}: two threads that wish to observe or modify the
    same object must coordinate.

  \item \textbf{Serialization ambiguity.}  When transmitting an
    object graph across a network, the serialization framework must
    decide whether aliased references should be preserved (deep
    structural copy) or collapsed (sharing-preserving copy).  This
    decision is invisible at the type level and frequently leads to
    subtle bugs in distributed systems~\cite{armstrong2007programming}.

  \item \textbf{Cache invalidation.}  Mutable state requires cache
    invalidation protocols---one of the ``two hard things in computer
    science''~\cite{fowler2003patterns}.  An immutable value, once
    cached, never becomes stale.

  \item \textbf{Testing difficulty.}  Objects with mutable state are
    difficult to test in isolation because their behaviour depends on
    the history of mutations.  Values, being immutable, have no
    history; they are their own specification.
\end{enumerate}

\subsection{Our Thesis}

We contend that treating \emph{values as primary}---making immutable
algebraic data types the default representation for all data, and
confining mutation to an explicitly tracked resource layer---yields a
language design that is:

\begin{itemize}
  \item \textbf{Simpler:} local reasoning is restored; every binding
    denotes a fixed value.
  \item \textbf{Safer:} data races are eliminated by construction for
    all pure data.
  \item \textbf{Composable:} values compose freely via functions,
    without the fragile coupling of shared mutable state.
  \item \textbf{Distributable:} values can be serialized, transmitted,
    cached, and replicated without ambiguity.
\end{itemize}

This paper develops this thesis in the context of \japl{}, a strict,
typed, effect-aware functional programming language.  \japl{}'s
design motto is: ``Pure by default, concurrent by design,
resource-safe by construction, distributed without apology.''  The
\emph{Values Are Primary} principle is the first of \japl{}'s seven
core principles, and it is foundational: the remaining six principles
build upon the guarantees that pervasive immutability provides.

\subsection{Contributions}

This paper makes the following contributions:

\begin{enumerate}
  \item A formal framework for value semantics grounded in category
    theory, modelling algebraic data types as initial algebras of
    polynomial endofunctors (\S\ref{sec:formal}).
  \item A denotational semantics for a value-centric core calculus
    that guarantees referential transparency
    (\S\ref{sec:formal:denotational}).
  \item A detailed description of \japl{}'s value model, including
    immutable algebraic data types, structural records, tagged unions,
    and pattern matching (\S\ref{sec:japl-value-model}).
  \item A type system design combining parametric polymorphism,
    bidirectional inference, row polymorphism, and opaque types
    (\S\ref{sec:type-system}).
  \item A systematic comparison with eight other languages
    (\S\ref{sec:comparison}).
  \item Implementation strategies for efficient persistent data
    structures (\S\ref{sec:implementation}).
  \item Benchmarks demonstrating that value semantics need not
    sacrifice performance (\S\ref{sec:evaluation}).
  \item An analysis of the implications of default immutability for
    concurrent and distributed programming
    (\S\ref{sec:concurrency}).
\end{enumerate}

\subsection{Paper Organisation}

Section~\ref{sec:background} surveys background and related work.
Section~\ref{sec:formal} develops the formal framework.
Section~\ref{sec:japl-value-model} presents \japl{}'s value model.
Section~\ref{sec:type-system} describes the type system.
Section~\ref{sec:comparison} compares with related languages.
Section~\ref{sec:implementation} discusses implementation strategies.
Section~\ref{sec:evaluation} presents benchmarks.
Section~\ref{sec:concurrency} analyses concurrency implications.
Section~\ref{sec:discussion} discusses trade-offs and the dual-layer
model.  Section~\ref{sec:conclusion} concludes.


% ══════════════════════════════════════════════════════════════════
\section{Background and Related Work}
\label{sec:background}
% ══════════════════════════════════════════════════════════════════

\subsection{The Value of Values}

Rich Hickey's influential talk \emph{The Value of Values}~\cite{hickey2012value}
articulated a distinction that practitioners had felt but rarely named:
the difference between \emph{values} (immutable, timeless facts) and
\emph{places} (mutable, time-varying locations).  Hickey argued that
conflating the two---the default in most mainstream languages---is the
root cause of much accidental complexity.  A \emph{value} is a piece
of information; it does not change.  The number $42$, the date March 15,
the set $\{1,2,3\}$---these are values.  A ``mutable integer variable''
is not a value but a \emph{place} that happens to contain an integer
at this moment in time.

Hickey's Clojure~\cite{hickey2008clojure} operationalised this insight
by making all core data structures---lists, vectors, hash maps,
sets---immutable and persistent.  The language provides controlled
mutation through \emph{atoms} (compare-and-swap references),
\emph{refs} (software transactional memory), and \emph{agents}
(asynchronous state updates), all of which are clearly delineated from
pure values.

\subsection{ML and the Algebraic Tradition}

The ML family~\cite{milner1978theory,milner1997definition} introduced
algebraic data types and pattern matching as core language features.
Standard ML~\cite{milner1997definition} and its descendants---OCaml~\cite{leroy2014ocaml},
F\#~\cite{syme2010expert}---demonstrate the power of sum types (tagged
unions) and product types (records/tuples) for modelling data.
However, ML languages generally permit unrestricted mutation via
\texttt{ref} cells, meaning that the value-centric idiom is a
\emph{convention} rather than a \emph{guarantee}.

\subsection{Haskell: Purity by Default}

Haskell~\cite{jones2003haskell,marlow2010haskell} takes the value-centric
approach to its logical extreme: all expressions denote values, and
side effects are confined to the \texttt{IO} monad~\cite{peyton1993imperative}.
This provides the strongest possible guarantees---referential transparency
holds for all non-\texttt{IO} code---but the monadic encoding of
effects introduces syntactic overhead and monad transformer stacks that
many practitioners find difficult~\cite{kiselyov2013extensible}.  Furthermore,
Haskell's lazy evaluation strategy means that ``values'' may in fact be
unevaluated thunks, complicating space usage reasoning~\cite{jones1999playing}.

\subsection{Erlang/OTP: Values in the Actor Model}

Erlang~\cite{armstrong2007programming,armstrong1996concurrent} combines
immutable values with the actor model.  All data in Erlang is
immutable; variables are single-assignment.  Processes communicate by
sending messages, which are deep-copied into the recipient's mailbox.
This design---immutable values plus process isolation---has proven
extraordinarily robust in telecommunications and distributed systems.
The main limitation is the lack of a static type system; Erlang is
dynamically typed, and the Dialyzer~\cite{lindahl2006practical}
provides only post-hoc analysis.  Elixir~\cite{valim2014elixir}
builds on the Erlang VM with improved syntax but retains the same
value model.  Gleam~\cite{gleam2024} adds static typing to the BEAM
ecosystem while preserving Erlang's value semantics.

\subsection{Rust: Ownership Without Garbage Collection}

Rust~\cite{matsakis2014rust,klabnik2019rust} takes a different approach
to the mutation problem.  Rather than prohibiting mutation, Rust
controls \emph{aliasing}: through the ownership and borrowing system,
Rust ensures that mutable references are unique (no aliasing) while
shared references are immutable (no mutation).  This ``aliasing XOR
mutability'' discipline~\cite{jung2017rustbelt} eliminates data races
without requiring immutability.  Values in Rust are moved by default
and copied only when explicitly requested via the \texttt{Copy} or
\texttt{Clone} traits.

\subsection{Persistent Data Structures}

Okasaki's seminal work~\cite{okasaki1999purely} demonstrated that
purely functional data structures can achieve asymptotically efficient
operations through techniques such as lazy evaluation, amortisation,
and structural decomposition.  Bagwell's hash-array mapped
tries~\cite{bagwell2001ideal} provided the practical basis for
Clojure's persistent vectors and hash maps, achieving near-constant-time
operations with $O(\log_{32} n)$ path-copying overhead.  Driscoll
et al.~\cite{driscoll1989making} formalised the theory of persistent
data structures, distinguishing partial from full persistence and
establishing lower bounds.

\subsection{Category-Theoretic Foundations}

The connection between algebraic data types and initial algebras in
category theory was established by Hagino~\cite{hagino1987category}
and further developed by Malcolm~\cite{malcolm1990algebraic} and
Meijer et al.~\cite{meijer1991functional}.  The insight is that a
data type definition such as $\texttt{List}(a) = \texttt{Nil} \mid
\texttt{Cons}(a, \texttt{List}(a))$ defines a functor $F(X) = 1 + a
\times X$, and the type \texttt{List}$(a)$ is the carrier of the
initial $F$-algebra.  This gives a principled account of structural
recursion (catamorphisms) and induction on data types.  Pierce's
\emph{Types and Programming Languages}~\cite{pierce2002types} and
Harper's \emph{Practical Foundations for Programming
Languages}~\cite{harper2016practical} provide comprehensive
treatments of the type-theoretic perspective.


% ══════════════════════════════════════════════════════════════════
\section{Formal Framework}
\label{sec:formal}
% ══════════════════════════════════════════════════════════════════

\subsection{Types as Objects of a Category}

We model the type system of our value-centric calculus using the
framework of \emph{Cartesian closed categories}
(CCCs)~\cite{lambek1986introduction,barr1990category}.

\begin{definition}[Category of Types]
\label{def:type-category}
Let $\Type$ be a category whose objects are types and whose morphisms
$f : A \to B$ are (total, terminating) programs that take a value of
type $A$ and produce a value of type $B$.  Composition is function
composition; the identity morphism at $A$ is the identity function
$\id_A : A \to A$.
\end{definition}

\begin{definition}[Cartesian Closed Structure]
\label{def:ccc}
$\Type$ is Cartesian closed if it has:
\begin{enumerate}
  \item A \emph{terminal object} $\mathbf{1}$ (the unit type).
  \item \emph{Binary products} $A \times B$ for all objects $A, B$
    (product types / tuples).
  \item \emph{Exponential objects} $B^A$ for all objects $A, B$
    (function types $A \to B$).
\end{enumerate}
The Cartesian closed structure ensures that products and function
types interact correctly: there is a natural bijection
$\Hom(A \times B, C) \cong \Hom(A, C^B)$ (currying).
\end{definition}

\begin{remark}
The Cartesian closed structure captures the essence of
\emph{value-centric} programming: products model data aggregation
(records, tuples), exponentials model computation (functions), and the
terminal object models the trivial value (unit).  No morphism has a
``side effect''; all information flows through inputs and outputs.
\end{remark}

\subsection{Algebraic Data Types as Initial Algebras}
\label{sec:formal:initial-algebras}

\begin{definition}[Polynomial Endofunctor]
\label{def:poly-functor}
A \emph{polynomial endofunctor} $F : \Type \to \Type$ is built from
the grammar:
\[
  F ::= \mathrm{Id} \mid K_A \mid F_1 + F_2 \mid F_1 \times F_2
\]
where $\mathrm{Id}$ is the identity functor, $K_A$ is the constant
functor at $A$, $+$ is the coproduct (sum), and $\times$ is the
product.
\end{definition}

\begin{definition}[$F$-Algebra]
\label{def:f-algebra}
An $F$-algebra is a pair $(A, \alpha)$ where $A$ is an object of
$\Type$ (the \emph{carrier}) and $\alpha : F(A) \to A$ is a morphism
(the \emph{structure map}).  A homomorphism of $F$-algebras
$(A, \alpha) \to (B, \beta)$ is a morphism $h : A \to B$ such that
$h \circ \alpha = \beta \circ F(h)$.
\end{definition}

\begin{definition}[Initial Algebra]
\label{def:initial-algebra}
An $F$-algebra $(\mu F, \mathrm{in})$ is \emph{initial} if, for every
$F$-algebra $(A, \alpha)$, there exists a unique homomorphism
$\sem{\alpha} : \mu F \to A$ satisfying:
\[
  \sem{\alpha} \circ \mathrm{in} = \alpha \circ F(\sem{\alpha})
\]
The morphism $\sem{\alpha}$ is called the \emph{catamorphism} (or
fold) induced by $\alpha$.
\end{definition}

\begin{theorem}[Lambek's Lemma~\cite{lambek1968fixpoint}]
\label{thm:lambek}
If $(\mu F, \mathrm{in})$ is an initial $F$-algebra, then
$\mathrm{in} : F(\mu F) \to \mu F$ is an isomorphism.  That is,
$\mu F \cong F(\mu F)$.
\end{theorem}

\begin{example}[Lists as an Initial Algebra]
\label{ex:list-initial}
The type $\texttt{List}(a)$ is the initial algebra of the functor
$F_a(X) = 1 + a \times X$.  The structure map is:
\[
  \mathrm{in} = [\texttt{Nil}, \texttt{Cons}] : 1 + a \times \texttt{List}(a) \to \texttt{List}(a)
\]
Given any algebra $(B, [z, f])$ where $z : 1 \to B$ and $f : a \times
B \to B$, the catamorphism $\sem{z, f} : \texttt{List}(a) \to B$ is
the unique function satisfying:
\[
  \sem{z, f}(\texttt{Nil}) = z(\ast) \qquad
  \sem{z, f}(\texttt{Cons}(x, xs)) = f(x, \sem{z, f}(xs))
\]
This is precisely the familiar \texttt{fold} operation.
\end{example}

\begin{example}[JAPL's Result Type]
\label{ex:result-initial}
In \japl{}, the \texttt{Result} type is defined as:
\begin{lstlisting}
type Result(a, e) =
  | Ok(a)
  | Err(e)
\end{lstlisting}
This is the initial algebra of the functor $F_{a,e}(X) = a + e$,
which is actually a constant functor (no recursive occurrences).  The
initial algebra is simply $a + e$ itself.  More precisely,
\texttt{Result} is a \emph{bifunctor} $\Type \times \Type \to \Type$
that sends $(a, e)$ to the coproduct $a + e$.
\end{example}

\subsection{Value Semantics: Formal Properties}
\label{sec:formal:value-semantics}

\begin{definition}[Value Semantics]
\label{def:value-semantics}
A language has \emph{value semantics} if, for every well-typed
expression $e$ of type $\tau$:
\begin{enumerate}
  \item \textbf{No observable mutation:} If $e$ evaluates to a value
    $v$, then $v$ cannot be observably changed by any operation in
    the language.  Formally, for all contexts $C[-]$:
    \[
      C[e] \Downarrow v \implies \forall C'[-].\; C'[e] \Downarrow v
    \]
  \item \textbf{Referential transparency:} In any expression, $e$
    can be replaced by its value $v$ without changing the meaning of
    the program.  Formally, if $e \Downarrow v$, then for all
    contexts $C[-]$:
    \[
      \sem{C[e]} = \sem{C[v]}
    \]
  \item \textbf{Structural equality:} Two values of the same
    \emph{first-order} type (products, sums, records, primitive types)
    are equal if and only if they have the same structure.  There is
    no notion of ``identity'' beyond structural content:
    \[
      v_1 = v_2 \iff \mathrm{struct}(v_1) = \mathrm{struct}(v_2)
    \]
    \emph{Function types} (exponentials $B^A$) are excluded from
    decidable structural equality: extensional equality of functions
    is undecidable in general.  In \japl{}, function values do not
    implement the \texttt{Eq} trait; attempting to compare two
    closures is a compile-time error.  Equality is restricted to ADTs,
    records, and primitive types where structural comparison is
    well-defined and decidable.
\end{enumerate}
\end{definition}

\begin{theorem}[Compositionality of Value Semantics]
\label{thm:compositionality}
In a language with value semantics, the denotation of any expression
is determined solely by the denotations of its subexpressions.
\end{theorem}

\begin{proof}
By structural induction on the syntax of expressions.  In the base
case, a literal denotes itself.  In the inductive case, consider an
expression $f(e_1, \ldots, e_n)$.  Since $f$ is a pure function
(no observable mutation), its output is determined by its inputs.
By the induction hypothesis, each $\sem{e_i}$ is determined by the
denotations of its subexpressions.  Therefore $\sem{f(e_1, \ldots,
e_n)} = \sem{f}(\sem{e_1}, \ldots, \sem{e_n})$ is fully determined.
\end{proof}

\subsection{Denotational Semantics of a Value-Centric Core Calculus}
\label{sec:formal:denotational}

We define a small core calculus $\lambda_V$ that captures the essence
of \japl{}'s value layer.

\subsubsection{Syntax}

\[
\begin{array}{rcl}
  \tau & ::= & \mathbf{1} \mid \alpha \mid \tau_1 \to \tau_2
    \mid \tau_1 \times \tau_2 \mid \tau_1 + \tau_2
    \mid \mu\alpha.\tau \mid \{l_1:\tau_1, \ldots, l_n:\tau_n\} \\[6pt]
  e & ::= & () \mid x \mid \lambda x:\tau.\; e \mid e_1\; e_2
    \mid (e_1, e_2) \mid \pi_1(e) \mid \pi_2(e) \\
    & \mid & \mathrm{inl}(e) \mid \mathrm{inr}(e) \mid
    \mathbf{case}\; e\; \{x.e_1, y.e_2\} \\
    & \mid & \mathrm{fold}(e) \mid \mathrm{unfold}(e) \mid
    \{l_1=e_1, \ldots, l_n=e_n\} \mid e.l
\end{array}
\]

\subsubsection{Typing Rules (Selected)}

\begin{gather}
  \frac{}{\Gamma \vdash () : \mathbf{1}} \quad (\text{T-Unit})
  \qquad
  \frac{x:\tau \in \Gamma}{\Gamma \vdash x : \tau} \quad (\text{T-Var})
  \\[8pt]
  \frac{\Gamma, x:\tau_1 \vdash e : \tau_2}
       {\Gamma \vdash \lambda x:\tau_1.\; e : \tau_1 \to \tau_2}
  \quad (\text{T-Abs})
  \qquad
  \frac{\Gamma \vdash e_1 : \tau_1 \to \tau_2 \quad \Gamma \vdash e_2 : \tau_1}
       {\Gamma \vdash e_1\; e_2 : \tau_2}
  \quad (\text{T-App})
  \\[8pt]
  \frac{\Gamma \vdash e : \tau[\mu\alpha.\tau/\alpha]}
       {\Gamma \vdash \mathrm{fold}(e) : \mu\alpha.\tau}
  \quad (\text{T-Fold})
  \qquad
  \frac{\Gamma \vdash e : \mu\alpha.\tau}
       {\Gamma \vdash \mathrm{unfold}(e) : \tau[\mu\alpha.\tau/\alpha]}
  \quad (\text{T-Unfold})
  \\[8pt]
  \frac{\Gamma \vdash e : \tau_1 + \tau_2 \quad
        \Gamma, x:\tau_1 \vdash e_1 : \tau \quad
        \Gamma, y:\tau_2 \vdash e_2 : \tau}
       {\Gamma \vdash \mathbf{case}\; e\; \{x.e_1,\, y.e_2\} : \tau}
  \quad (\text{T-Case})
  \\[8pt]
  \frac{\Gamma \vdash e : \tau_1}
       {\Gamma \vdash \mathrm{inl}(e) : \tau_1 + \tau_2}
  \quad (\text{T-Inl})
  \qquad
  \frac{\Gamma \vdash e : \tau_2}
       {\Gamma \vdash \mathrm{inr}(e) : \tau_1 + \tau_2}
  \quad (\text{T-Inr})
\end{gather}

\subsubsection{Denotational Semantics}

We interpret types as sets and terms as set-theoretic functions.

\begin{definition}[Type Interpretation]
\label{def:type-interp}
\begin{align}
  \sem{\mathbf{1}} &= \{\ast\} \\
  \sem{\tau_1 \to \tau_2} &= \sem{\tau_2}^{\sem{\tau_1}}
    \quad\text{(function space)} \\
  \sem{\tau_1 \times \tau_2} &= \sem{\tau_1} \times \sem{\tau_2}
    \quad\text{(Cartesian product)} \\
  \sem{\tau_1 + \tau_2} &= \sem{\tau_1} \uplus \sem{\tau_2}
    \quad\text{(disjoint union)} \\
  \sem{\mu\alpha.\tau} &= \mathrm{Fix}(\lambda X.\, \sem{\tau}[X/\alpha])
    \quad\text{(least fixpoint)} \\
  \sem{\{l_1:\tau_1, \ldots, l_n:\tau_n\}} &=
    \sem{\tau_1} \times \cdots \times \sem{\tau_n}
    \quad\text{(labelled product)}
\end{align}
\end{definition}

\begin{definition}[Term Interpretation]
\label{def:term-interp}
Given an environment $\rho : \mathrm{Var} \to \bigcup_\tau \sem{\tau}$:
\begin{align}
  \sem{()}^\rho &= \ast \\
  \sem{x}^\rho &= \rho(x) \\
  \sem{\lambda x:\tau.\; e}^\rho &=
    \lambda v \in \sem{\tau}.\; \sem{e}^{\rho[x \mapsto v]} \\
  \sem{e_1\; e_2}^\rho &= \sem{e_1}^\rho(\sem{e_2}^\rho) \\
  \sem{(e_1, e_2)}^\rho &= (\sem{e_1}^\rho, \sem{e_2}^\rho) \\
  \sem{\mathrm{inl}(e)}^\rho &= \mathrm{inl}(\sem{e}^\rho) \\
  \sem{\mathbf{case}\; e\; \{x.e_1, y.e_2\}}^\rho &=
    \begin{cases}
      \sem{e_1}^{\rho[x \mapsto v]} & \text{if } \sem{e}^\rho = \mathrm{inl}(v) \\
      \sem{e_2}^{\rho[y \mapsto v]} & \text{if } \sem{e}^\rho = \mathrm{inr}(v)
    \end{cases}
\end{align}
\end{definition}

\begin{theorem}[Adequacy]
\label{thm:adequacy}
The denotational semantics is adequate with respect to the operational
semantics: if $e \Downarrow v$ (big-step), then $\sem{e} = \sem{v}$.
Conversely, if $\sem{e} \neq \bot$, then $e$ terminates.
\end{theorem}

\begin{proof}[Proof sketch]
Adequacy follows from the standard technique of defining a logical
relation $\mathcal{R}_\tau \subseteq \sem{\tau} \times
\mathrm{Terms}_\tau$ between denotations and terms, and showing by
structural induction that (1) if $e \Downarrow v$ then $(\sem{e},
v) \in \mathcal{R}_\tau$, and (2) if $(d, e) \in \mathcal{R}_\tau$
and $d \neq \bot$, then $e \Downarrow v$ for some $v$ with
$\sem{v} = d$.  The key insight is that in a value-centric calculus
without mutation, the logical relation can be defined without
step-indexing, since there are no store-dependent types to account
for.  This follows the approach of Plotkin~\cite{plotkin1977lcf}
as adapted by Harper~\cite{harper2016practical}.
\end{proof}

\begin{corollary}[Referential Transparency]
\label{cor:ref-trans}
For all expressions $e$ and contexts $C[-]$ in $\lambda_V$:
if $e \Downarrow v$, then $\sem{C[e]} = \sem{C[v]}$.
\end{corollary}

\begin{proof}
Immediate from compositionality (Theorem~\ref{thm:compositionality})
and adequacy (Theorem~\ref{thm:adequacy}).
\end{proof}

\subsection{The Yoneda Perspective on Values}
\label{sec:formal:yoneda}

The Yoneda lemma~\cite{maclane1998categories} provides an elegant
characterisation of values in our framework.

\begin{theorem}[Yoneda Lemma]
For any locally small category $\mathcal{C}$, object $A$, and functor
$F : \mathcal{C}^{\mathrm{op}} \to \Set$:
\[
  \mathrm{Nat}(\Hom(-, A), F) \cong F(A)
\]
naturally in $A$.
\end{theorem}

In the category $\Type$, the Yoneda lemma tells us that a value
$v : A$ is completely determined by the collection of all morphisms
\emph{out of} $A$---that is, by the functions that can act on $v$.
This is the formal content of the principle that ``a value \emph{is}
what you can do with it.''  Two values that are indistinguishable by
all observers are, by Yoneda, identical.  This is precisely the
structural equality property of Definition~\ref{def:value-semantics}.

\begin{remark}
In an identity-centric language, two objects can be observationally
equivalent yet have distinct identities (distinct heap locations).
The Yoneda lemma shows that in a value-centric language, this
situation cannot arise: values have no identity beyond their
observable behaviour.
\end{remark}

\subsubsection{Yoneda in JAPL's Design}

The Yoneda perspective directly informs three concrete aspects of
\japl{}'s implementation:

\begin{enumerate}
  \item \textbf{Compiler-derived equality.}  The \texttt{deriving(Eq)}
    mechanism generates structural equality tests by recursively
    comparing fields.  This is justified by Yoneda: two values that
    agree on all observations (including every projection and pattern
    match) must be structurally identical, so the compiler need only
    check structure.

  \item \textbf{Content-addressable hashing.}  Because a value \emph{is}
    its observable behaviour, \japl{} can safely derive \texttt{Hash}
    implementations that hash the structure of a value.  Structural
    identity implies hash identity---the Yoneda condition guarantees
    this is sound, and it underpins the HAMT-based collections
    (\S\ref{sec:implementation:hamt}).

  \item \textbf{Opaque-type abstraction.}  An opaque type
    (\S\ref{sec:type-system:opaque}) restricts the set of morphisms
    available to external modules, thereby creating a smaller
    representable functor.  Inside the defining module, the full
    Yoneda embedding applies; outside, clients see only the exported
    interface.  This is the operational content of the Yoneda lemma
    applied to module boundaries.
\end{enumerate}


% ══════════════════════════════════════════════════════════════════
\section{JAPL's Value Model}
\label{sec:japl-value-model}
% ══════════════════════════════════════════════════════════════════

\subsection{Immutable by Default}

In \japl{}, every binding introduces a value.  There is no
\keyword{let mut} or \keyword{var}; bindings are irrevocable.

\begin{lstlisting}[caption={Basic value bindings in \japl{}.}]
-- Values are the default. No `let mut`, no `var`.
let name = "JAPL"
let point = { x = 3.0, y = 4.0 }
let items = [1, 2, 3, 4, 5]

-- Transformation produces new values, never mutates.
let shifted = { point | x = point.x + 1.0 }
let doubled = List.map(items, fn x -> x * 2)
\end{lstlisting}

The record update syntax \texttt{\{ point | x = ... \}} constructs a
new record that shares structure with the original.  The original
\texttt{point} is unchanged and remains accessible.

\subsection{Algebraic Data Types}
\label{sec:japl:adt}

\japl{} provides algebraic data types (ADTs) as the primary mechanism
for defining structured data.  ADTs encompass both \emph{product types}
(records with named fields) and \emph{sum types} (tagged unions with
variants).

\subsubsection{Product Types: Structural Records}

\begin{lstlisting}[caption={Record types in \japl{}.}]
type User = {
  id: UserId,
  name: String,
  email: String
}

type Point = { x: Float, y: Float }

-- Records are created with `=` syntax
let alice = { id = UserId(1), name = "Alice",
              email = "alice@example.com" }
\end{lstlisting}

Records in \japl{} are \emph{structurally typed}: two record types with
the same field names and types are compatible, regardless of whether
they were defined with the same \keyword{type} declaration.  This is
formalised via row polymorphism (\S\ref{sec:type-system:row}).

\subsubsection{Sum Types: Tagged Unions}

\begin{lstlisting}[caption={Sum types in \japl{}.}]
type Result(a, e) =
  | Ok(a)
  | Err(e)

type Shape =
  | Circle(Float)
  | Rectangle(Float, Float)
  | Triangle(Float, Float, Float)

type Option(a) =
  | Some(a)
  | None
\end{lstlisting}

Sum types are \emph{closed}: the set of variants is fixed at definition
time.  This enables exhaustive pattern matching, which the compiler
enforces.

\subsubsection{Pattern Matching}

Pattern matching is the primary mechanism for deconstructing values:

\begin{lstlisting}[caption={Pattern matching in \japl{}.}]
fn area(shape: Shape) -> Float =
  match shape with
  | Circle(r) -> 3.14159 * r * r
  | Rectangle(w, h) -> w * h
  | Triangle(a, b, c) ->
      let s = (a + b + c) / 2.0
      Float.sqrt(s * (s - a) * (s - b) * (s - c))

fn user_label(user: User) -> String =
  user.name <> " <" <> user.email <> ">"
\end{lstlisting}

The compiler checks pattern matches for exhaustiveness and redundancy.
A missing case is a compile-time error; an unreachable case is a
warning.  This is essential for maintainability: adding a variant to
a sum type causes all incomplete matches to be flagged.

\subsection{The Value/Resource Split}
\label{sec:japl:split}

Not all data can be treated as a pure value.  File handles, network
sockets, GPU buffers, and database connections are \emph{resources}:
entities with identity, lifecycle, and finality.  A socket that has
been closed cannot be ``used again''; a file handle must be released
exactly once.

\japl{} addresses this through a clean two-layer architecture:

\begin{enumerate}
  \item \textbf{Pure Layer:} Immutable values managed by a garbage
    collector.  This is where algebraic data types, records, strings,
    lists, maps, and closures live.  Values are freely shareable
    across processes.

  \item \textbf{Resource Layer:} Mutable resources managed by linear
    types and ownership tracking.  Resources have single ownership,
    support borrowing, and are deterministically freed when their
    owner goes out of scope.
\end{enumerate}

\begin{lstlisting}[caption={The two-layer architecture.}]
-- Pure layer: GC-managed, immutable, freely shareable
let data = [1, 2, 3, 4, 5]
let copy = data  -- sharing is fine; data is immutable

-- Resource layer: ownership-tracked, must be consumed
fn process_file(path: String) -> Result(String, IoError) with Io =
  use file = File.open(path, Read)?
  let contents = File.read_all(file)?
  File.close(file)  -- consumed here; forgetting = compile error
  Ok(contents)
\end{lstlisting}

The \keyword{use} keyword introduces a linear binding: the variable
\texttt{file} must be consumed exactly once.  The compiler tracks
ownership and rejects programs that leak resources or use them after
transfer.

\begin{definition}[Value/Resource Separation]
\label{def:value-resource}
A type $\tau$ in \japl{} is classified as either a \emph{value type}
or a \emph{resource type}:
\begin{itemize}
  \item \textbf{Value types} satisfy the structural rules of
    weakening and contraction: values can be freely discarded or
    duplicated.
  \item \textbf{Resource types} are \emph{linear}: they must be used
    exactly once.  They do not admit weakening (cannot be discarded
    without explicit release) or contraction (cannot be duplicated).
\end{itemize}
The compiler statically enforces this classification, ensuring that
resource types are never treated as values and vice versa.

Crucially, \textbf{value types cannot contain resource types as
fields}.  A record or ADT in the pure layer may not hold an owned
file handle or socket.  This constraint is analogous to requiring
values to be \texttt{Send + Sync} in Rust's terminology: anything
that lives in the GC-managed, freely-copyable value layer must
itself be freely copyable.  If a value could embed an owned resource,
then duplicating the value (which the GC may do implicitly via
structural sharing) would violate the resource's linearity
invariant.  When a function needs to operate on both values and
resources, the resource is passed as a separate linearly-typed
parameter, not embedded inside a value.
\end{definition}

\subsection{Persistent Data Structure Strategies}
\label{sec:japl:persistent}

Since values are immutable, ``updating'' a data structure means
constructing a new version.  Naive deep copying would be prohibitively
expensive.  \japl{} employs \emph{persistent data structures} with
\emph{structural sharing}: the new version shares most of its structure
with the old version, and only the changed parts are allocated anew.

\begin{lstlisting}[caption={Structural sharing in record updates.}]
let user = { id = UserId(1), name = "Alice",
             email = "alice@example.com" }

-- Only the `email` field is newly allocated; `id` and `name`
-- are shared with the original.
let updated = { user | email = "alice@newdomain.com" }
\end{lstlisting}

For larger data structures (maps, sets, vectors), \japl{} uses
hash-array mapped tries (HAMTs) as described in
\S\ref{sec:implementation:hamt}.  The key insight is that an ``update''
to a map with $n$ entries requires only $O(\log_{32} n)$ new
allocations, sharing the vast majority of the tree with the original.


% ══════════════════════════════════════════════════════════════════
\section{Type System for Values}
\label{sec:type-system}
% ══════════════════════════════════════════════════════════════════

\subsection{Parametric Polymorphism}
\label{sec:type-system:poly}

\japl{} supports parametric polymorphism (generics) in the tradition
of System F~\cite{girard1972interpretation,reynolds1974towards}, but
with prenex quantification (quantifiers at the outermost level only)
for decidable type inference.

\begin{lstlisting}[caption={Parametric polymorphism in \japl{}.}]
fn map(list: List(a), f: fn(a) -> b) -> List(b) =
  match list with
  | [] -> []
  | [x, ..rest] -> [f(x), ..map(rest, f)]

fn compose(f: fn(b) -> c, g: fn(a) -> b) -> fn(a) -> c =
  fn x -> f(g(x))

fn identity(x: a) -> a = x
\end{lstlisting}

The type variables \texttt{a}, \texttt{b}, \texttt{c} are implicitly
universally quantified.  The compiler infers the most general type
for each function.

\begin{definition}[Parametricity]
\label{def:parametricity}
A polymorphic function $f : \forall \alpha.\, \tau(\alpha)$ satisfies
the \emph{parametricity} (or ``free theorem'')
condition~\cite{wadler1989theorems}: for any types $A, B$ and
function $g : A \to B$:
\[
  \sem{\tau}(g)(\sem{f}_A) = \sem{f}_B
\]
where $\sem{\tau}(g)$ is the action of $\tau$ on morphisms (viewing
$\tau$ as a functor).
\end{definition}

Parametricity ensures that polymorphic functions cannot ``peek'' at the
representation of their type parameters.  This gives us
\emph{free theorems}: for example, any function $f : \forall \alpha.\,
\texttt{List}(\alpha) \to \texttt{List}(\alpha)$ must commute with
\texttt{map}:
\[
  \texttt{map}(g, f(xs)) = f(\texttt{map}(g, xs))
\]
for all $g$ and $xs$.  This is a powerful reasoning tool enabled by
value semantics: in a language with mutation, a polymorphic function
could observe the representation of $\alpha$ through side effects,
violating parametricity.

\subsection{Type Inference: Local Bidirectional Checking}
\label{sec:type-system:inference}

\japl{} uses a bidirectional type checking
algorithm~\cite{pierce2000local,dunfield2021bidirectional} that
combines two modes:

\begin{enumerate}
  \item \textbf{Checking mode ($\Gamma \vdash e \Leftarrow \tau$):}
    Given a term $e$ and an expected type $\tau$, verify that $e$
    has type $\tau$.
  \item \textbf{Synthesis mode ($\Gamma \vdash e \Rightarrow \tau$):}
    Given a term $e$, infer its type $\tau$.
\end{enumerate}

\begin{gather}
  \frac{\Gamma \vdash e \Rightarrow \tau}
       {\Gamma \vdash e \Leftarrow \tau}
  \quad (\text{Sub})
  \\[8pt]
  \frac{\Gamma, x:\tau_1 \vdash e \Leftarrow \tau_2}
       {\Gamma \vdash \lambda x:\tau_1.\, e \Rightarrow \tau_1 \to \tau_2}
  \quad (\text{Abs-Synth})
  \\[8pt]
  \frac{\Gamma \vdash e_1 \Rightarrow \tau_1 \to \tau_2 \quad
        \Gamma \vdash e_2 \Leftarrow \tau_1}
       {\Gamma \vdash e_1\; e_2 \Rightarrow \tau_2}
  \quad (\text{App-Synth})
\end{gather}

Top-level function signatures are required at module boundaries, which
serves as both documentation and a firewall for type inference: the
compiler need not perform global inference.

\begin{lstlisting}[caption={Type inference in practice.}]
-- Signature required at module boundary
fn process(items: List(Item)) -> Summary with Io =
  -- Types inferred within the body
  let totals = List.map(items, fn item -> item.price * item.quantity)
  let sum = List.fold(totals, 0, fn acc, t -> acc + t)
  { item_count = List.length(items), total = sum }
\end{lstlisting}

\subsection{Row Polymorphism for Extensible Records}
\label{sec:type-system:row}

\japl{} supports row polymorphism~\cite{wand1991type,remy1994type},
allowing functions to operate on records with a minimum set of required
fields while remaining agnostic to additional fields.

\begin{definition}[Row Types]
\label{def:row-types}
A \emph{row} is a partial function from labels to types:
\[
  \rho : \mathrm{Label} \rightharpoonup \Type
\]
A record type $\{l_1:\tau_1, \ldots, l_n:\tau_n \mid \rho\}$ specifies
$n$ known fields and a \emph{row variable} $\rho$ representing
additional unknown fields.  Row unification equates rows modulo field
ordering.
\end{definition}

\begin{lstlisting}[caption={Row polymorphism in \japl{}.}]
-- Works on ANY record with a `name: String` field
fn greet(person: { name: String | r }) -> String =
  "Hello, " <> person.name <> "!"

-- All of these work:
let _ = greet({ name = "Alice" })
let _ = greet({ name = "Bob", age = 30 })
let _ = greet({ name = "Charlie", role = Admin })
\end{lstlisting}

Row polymorphism gives \japl{} a form of structural subtyping that is
strictly more principled than the duck typing of dynamically typed
languages or the interface-based structural typing of Go.  Every row-polymorphic
function has a precise type that the compiler can check and infer.

\begin{theorem}[Principal Types for Row Polymorphism]
\label{thm:row-principal}
Under R\'{e}my's row type discipline~\cite{remy1994type} with
equi-recursive row types, every well-typed expression has a principal
type, and type inference is decidable in polynomial time.
\end{theorem}

\subsection{Opaque and Newtype Wrappers}
\label{sec:type-system:opaque}

\japl{} supports \keyword{opaque} types for information hiding and
newtype wrappers for domain modelling:

\begin{lstlisting}[caption={Opaque types and newtypes.}]
-- Opaque type: implementation hidden from external modules
module Map

opaque type Map(k, v)

fn empty() -> Map(k, v) = ...
fn insert(map: Map(k, v), key: k, value: v) -> Map(k, v)
    where Ord(k) = ...

-- Newtype: zero-cost wrapper for type safety
type UserId = UserId(Int)
type Email = Email(String)

-- These are different types; cannot be mixed up
fn find_user(id: UserId) -> Option(User) = ...
\end{lstlisting}

Opaque types enforce abstraction boundaries: the internal
representation of \texttt{Map(k, v)} is visible within the
\texttt{Map} module but hidden from clients.  This enables changing
the implementation (e.g., from a balanced tree to a HAMT) without
affecting client code.

\subsection{Traits as Type Classes}
\label{sec:type-system:traits}

\japl{} uses traits (type classes) for ad-hoc polymorphism:

\begin{lstlisting}[caption={Traits in \japl{}.}]
trait Eq(a) =
  fn eq(x: a, y: a) -> Bool

trait Ord(a) where Eq(a) =
  fn compare(x: a, y: a) -> Ordering

trait Show(a) =
  fn show(value: a) -> String

trait Serialize(a) =
  fn serialize(value: a) -> Bytes
  fn deserialize(data: Bytes) -> Result(a, SerializeError)

-- Deriving: compiler generates implementations
type Point deriving(Eq, Ord, Show, Serialize) =
  { x: Float, y: Float }
\end{lstlisting}

Trait resolution uses Haskell-style dictionary passing, which interacts
well with value semantics: dictionaries are themselves values and can
be freely shared across processes.


% ══════════════════════════════════════════════════════════════════
\section{Comparison of Value Models}
\label{sec:comparison}
% ══════════════════════════════════════════════════════════════════

We compare the value models of eight languages along seven dimensions.
Table~\ref{tab:comparison} provides a summary; the following
subsections discuss each language in detail.

\begin{table}[H]
\centering
\small
\caption{Comparison of value models across languages.}
\label{tab:comparison}
\begin{tabularx}{\textwidth}{l *{4}{>{\centering\arraybackslash}X}}
\toprule
\textbf{Language} & \textbf{Default Immutable} & \textbf{Static Types}
  & \textbf{Persistent DS} & \textbf{Value Sharing} \\
\midrule
\textbf{\japl{}}  & Yes & Yes (inferred) & Yes (HAMT) & Across processes \\
Haskell           & Yes & Yes (inferred) & Yes (lazy)  & GHC RTS sharing \\
Erlang            & Yes & No (Dialyzer)  & Copy-based  & Deep copy to proc \\
Elixir            & Yes & No (Dialyzer)  & Copy-based  & Deep copy to proc \\
Gleam             & Yes & Yes (inferred) & Copy-based  & Deep copy to proc \\
Clojure           & Yes & No (Spec)      & Yes (HAMT)  & STM/Atom refs \\
Rust              & No  & Yes (inferred) & No (default)& Move semantics \\
OCaml             & No  & Yes (inferred) & Convention  & GC-managed refs \\
Go                & No  & Yes (declared) & No          & Goroutine + chan \\
\bottomrule
\end{tabularx}
\end{table}

\begin{table}[H]
\centering
\small
\caption{Comparison of value models (continued).}
\label{tab:comparison2}
\begin{tabularx}{\textwidth}{l *{3}{>{\centering\arraybackslash}X}}
\toprule
\textbf{Language} & \textbf{Pattern Matching}
  & \textbf{ADTs} & \textbf{Resource Mgmt} \\
\midrule
\textbf{\japl{}}  & Exhaustive & Sum + Product & Linear types \\
Haskell           & Exhaustive & Sum + Product & GC (no linear) \\
Erlang            & Exhaustive & Tuples/Maps   & GC per-process \\
Elixir            & Exhaustive & Structs/Maps  & GC per-process \\
Gleam             & Exhaustive & Sum + Product & GC (BEAM) \\
Clojure           & Destructure & Protocols    & GC + try-finally \\
Rust              & Exhaustive & Sum + Product & Ownership \\
OCaml             & Exhaustive & Sum + Product & GC (no linear) \\
Go                & Switch only & Structs only & defer \\
\bottomrule
\end{tabularx}
\end{table}

\subsection{Haskell}

Haskell is the closest language to \japl{}'s value model in terms of
purity guarantees.  All expressions in Haskell are referentially
transparent outside the \texttt{IO} monad.  However, three differences
are significant:

\begin{enumerate}
  \item \textbf{Laziness:} Haskell's lazy evaluation means that a
    ``value'' may actually be an unevaluated thunk occupying
    unpredictable amounts of memory~\cite{jones1999playing}.  \japl{}
    is strict, making space usage predictable.

  \item \textbf{No resource layer:} Haskell lacks linear types in the
    core language (GHC extensions notwithstanding~\cite{bernardy2018linear}).
    Resources must be managed through bracket patterns or the
    \texttt{ResourceT} monad transformer, which provides weaker
    guarantees than \japl{}'s linear types.

  \item \textbf{Monadic syntax:} Haskell's do-notation, while powerful,
    introduces a syntactic barrier between pure and effectful code.
    \japl{}'s effect annotations (\keyword{with} clauses) are
    lighter-weight and do not require monadic binding operators.
\end{enumerate}

\subsection{Rust}

Rust's approach is complementary to \japl{}'s.  Where \japl{} achieves
safety through \emph{default immutability} (values cannot change, so
sharing is safe), Rust achieves safety through \emph{controlled
aliasing} (mutable references cannot be aliased)~\cite{jung2017rustbelt}.

\begin{enumerate}
  \item \textbf{Values are not the default:} In Rust, \texttt{let}
    bindings are immutable but \texttt{let mut} is pervasive.  Data
    structures are mutable by default, and persistent data structures
    are not part of the standard library.

  \item \textbf{Move semantics:} Rust's move semantics mean that
    ``sharing'' a value between two contexts requires explicit cloning.
    In \japl{}, pure values are freely shareable because they are
    immutable.

  \item \textbf{No garbage collection:} Rust's ownership model
    provides deterministic deallocation without GC.  \japl{} uses GC
    for immutable values (which is efficient because immutable data
    needs no write barriers) and ownership for resources.
\end{enumerate}

\subsection{Erlang and Elixir}

Erlang is the most direct ancestor of \japl{}'s process model.  Both
languages share the fundamental insight that immutable values combined
with process isolation eliminate data races.

\begin{enumerate}
  \item \textbf{Dynamic typing:} Erlang and Elixir are dynamically
    typed.  Type errors are caught at runtime, not compile time.
    \japl{} provides full static typing with inference.

  \item \textbf{Deep copying:} On the BEAM VM, messages sent between
    processes are deep-copied into the receiver's heap (with
    exceptions for large binaries stored in a shared heap).  \japl{}'s
    design allows zero-copy message passing for immutable values,
    since they cannot be mutated by the sender after sending.

  \item \textbf{No ADTs:} Erlang uses tuples and atoms for tagged data,
    which lacks the compiler-checked exhaustiveness of algebraic data
    types.  Gleam~\cite{gleam2024} addresses this within the BEAM
    ecosystem.
\end{enumerate}

\subsection{Gleam}

Gleam~\cite{gleam2024} occupies a position very close to \japl{} in
the design space: it is a statically typed, strict, functional language
targeting the BEAM VM with algebraic data types and exhaustive pattern
matching.

\begin{enumerate}
  \item \textbf{No effect system:} Gleam does not track effects in
    types.  Any function can perform I/O without annotation.  \japl{}
    tracks effects explicitly.

  \item \textbf{No resource layer:} Gleam relies on the BEAM's
    per-process GC for all memory management, including resources.
    \japl{} provides linear types for deterministic resource cleanup.

  \item \textbf{BEAM constraints:} Gleam inherits the BEAM's
    deep-copy message passing and runtime characteristics.  \japl{}
    targets a custom runtime that can exploit immutability for
    zero-copy sharing.
\end{enumerate}

\subsection{OCaml}

OCaml~\cite{leroy2014ocaml} is a strict, statically typed functional
language with algebraic data types---superficially similar to \japl{}.
The key difference is that OCaml permits unrestricted mutation through
\texttt{ref} cells and mutable record fields.  This means that OCaml
values are not necessarily values in the sense of
Definition~\ref{def:value-semantics}: a record may contain a
\texttt{ref} field that can be mutated.

\subsection{Clojure}

Clojure~\cite{hickey2008clojure} is the language that most explicitly
articulates the ``values are primary'' philosophy, through Hickey's
talks and writings~\cite{hickey2012value}.  Clojure's persistent data
structures (vectors, hash maps, sets) based on
HAMTs~\cite{bagwell2001ideal} are the practical gold standard.

However, Clojure is dynamically typed and runs on the JVM, which
means: (1) type errors are caught at runtime; (2) all values are
heap-allocated objects subject to JVM GC pauses; (3) the concurrency
model (STM, atoms, agents) is shared-memory rather than process-based.

\subsection{Go}

Go~\cite{donovan2015go} represents the opposite end of the spectrum.
Go values are mutable by default; there is no \texttt{const}
qualifier for compound types.  Go's concurrency model
(goroutines + channels) provides memory safety through the mantra
``share memory by communicating,'' but the language does not enforce
this---nothing prevents goroutines from sharing mutable pointers.


% ══════════════════════════════════════════════════════════════════
\section{Implementation Strategies}
\label{sec:implementation}
% ══════════════════════════════════════════════════════════════════

A common objection to value semantics is that immutability is
expensive: every ``update'' copies the entire data structure.  This
section demonstrates that with appropriate implementation techniques,
value semantics can achieve performance competitive with mutable
approaches.

\subsection{Hash-Array Mapped Tries (HAMTs)}
\label{sec:implementation:hamt}

The hash-array mapped trie~\cite{bagwell2001ideal} is the workhorse
data structure for persistent collections in value-centric languages.
A HAMT is a wide, shallow trie indexed by hash values, with a
branching factor of 32 (using 5 bits of the hash at each level).

\begin{definition}[HAMT Structure]
\label{def:hamt}
A HAMT node contains:
\begin{itemize}
  \item A 32-bit \emph{bitmap} indicating which children are present.
  \item A compact array of children, with $\mathrm{popcount}(\text{bitmap})$
    entries.
\end{itemize}
For a collection of $n$ elements, the trie has depth at most
$\lceil \log_{32} n \rceil = \lceil \frac{\log n}{5} \rceil$.
\end{definition}

\begin{theorem}[HAMT Complexity]
\label{thm:hamt-complexity}
For a HAMT with $n$ entries:
\begin{itemize}
  \item \textbf{Lookup:} $O(\log_{32} n)$ time (effectively $O(1)$
    for practical sizes).
  \item \textbf{Insert/Update:} $O(\log_{32} n)$ time and space
    (path copying).
  \item \textbf{Structural sharing:} An update shares
    $O(n - \log_{32} n)$ nodes with the original.
\end{itemize}
\end{theorem}

\begin{proof}
Lookup traverses from the root to a leaf, consuming 5 bits of the
32-bit hash at each level, giving depth $\leq 7$ for 32-bit hashes.
Insertion path-copies the $O(\log_{32} n)$ nodes on the path from
root to the insertion point, sharing all other nodes.
\end{proof}

In practice, with a branching factor of 32 and hash space of
$2^{32}$, the maximum depth is 7.  This means that inserting into a
map of a million entries requires copying at most 7 nodes (each of
size $\leq 32$ pointers)---a total of at most 224 words of allocation.

\subsection{Structural Sharing}
\label{sec:implementation:sharing}

The key principle enabling efficient persistent data structures is
\emph{structural sharing}: when a data structure is ``updated,'' the
new version shares all unchanged parts with the old version.

\begin{figure}[h]
\centering
\begin{verbatim}
  Original:        After insert("d", 4):

      [a b c]            [a b c d]
     / | \              / | \ \
   v1 v2 v3          v1 v2 v3 v4
                      ^^  ^^  ^^
                      shared with original
\end{verbatim}
\caption{Structural sharing in a persistent vector.  The original
  vector and the updated vector share the leaf nodes for elements
  \texttt{a}, \texttt{b}, \texttt{c}.}
\label{fig:structural-sharing}
\end{figure}

For records, structural sharing is even simpler: a record update
\texttt{\{ user | email = new\_email \}} allocates a new record
header pointing to the shared \texttt{id} and \texttt{name} fields
and the new \texttt{email} field.

\subsection{Copy-on-Write for Small Values}
\label{sec:implementation:cow}

For small values (records with few fields, short lists), full copying
can be more efficient than persistent data structures due to cache
locality.  \japl{}'s compiler employs a threshold-based strategy:

\begin{itemize}
  \item \textbf{Small records} ($\leq 8$ fields): copied in full on
    update.  The copy fits in one or two cache lines.
  \item \textbf{Large records} ($> 8$ fields): use pointer-based
    structural sharing.
  \item \textbf{Small lists} ($\leq 32$ elements): represented as
    flat arrays, copied on update.
  \item \textbf{Large lists}: represented as RRB-trees (relaxed
    radix-balanced trees)~\cite{stucki2015rrb} with structural
    sharing.
\end{itemize}

\subsection{Packed Tagged Unions}
\label{sec:implementation:tagged-unions}

Sum types are represented as tagged unions with compiler-optimised
layouts:

\begin{lstlisting}[caption={Packed representation for small unions.}]
type Vec3 = packed { x: Float32, y: Float32, z: Float32 }

type Particle = packed
  | Active(Vec3, Vec3, Float32)   -- position, velocity, mass
  | Inactive
\end{lstlisting}

\begin{definition}[Tagged Union Layout]
\label{def:tagged-layout}
A sum type with variants $C_1(\tau_{1,1}, \ldots, \tau_{1,k_1}),
\ldots, C_n(\tau_{n,1}, \ldots, \tau_{n,k_n})$ is laid out as:
\[
  \underbrace{\text{tag}}_{\lceil \log_2 n \rceil \text{ bits}}
  \;\Big|\;
  \underbrace{\text{payload}}_{\max_i \sum_j |\tau_{i,j}| \text{ bits}}
\]
The compiler chooses the smallest tag size that can distinguish all
variants and uses the maximum payload size across all variants.
\end{definition}

Additional optimisations include:

\begin{enumerate}
  \item \textbf{Null pointer optimisation:} For types like
    \texttt{Option(ref T)}, the \texttt{None} variant is represented
    as a null pointer, eliminating the tag entirely.
  \item \textbf{Tag-in-pointer:} When values are heap-allocated, the
    tag can be stored in the low bits of the pointer (which are
    otherwise unused due to alignment).
  \item \textbf{Unboxing:} Small sum types (e.g., \texttt{Bool},
    \texttt{Ordering}) are unboxed---stored directly in registers or
    on the stack, never heap-allocated.
\end{enumerate}

\subsection{Cache-Friendly Layouts}
\label{sec:implementation:cache}

\japl{} provides the \keyword{packed} modifier for types that should
be laid out contiguously in memory:

\begin{lstlisting}[caption={Cache-friendly layout with \texttt{packed}.}]
-- Contiguous array of Vec3 values (12 bytes each, no pointers)
type Vec3 = packed { x: Float32, y: Float32, z: Float32 }

-- An array of 1000 Vec3 values occupies 12KB of contiguous memory
let positions: Array(Vec3) = Array.create(1000, Vec3(0.0, 0.0, 0.0))
\end{lstlisting}

The \keyword{packed} modifier instructs the compiler to use a flat,
unboxed representation.  Fields are stored contiguously with no
pointer indirection.  This is critical for numerical and
data-intensive workloads where cache performance dominates.

\subsection{Uniqueness Analysis}
\label{sec:implementation:uniqueness}

When the compiler can determine that a value has a unique reference
(no other part of the program holds a reference to it), it can perform
the update \emph{in place}, converting a logical copy into a physical
mutation.  This is a form of \emph{uniqueness typing}~\cite{barendsen1996uniqueness}
applied as an optimisation.

\begin{lstlisting}[caption={Compiler-optimised in-place update.}]
fn build_list(n: Int) -> List(Int) =
  let rec go(i, acc) =
    if i == 0 then acc
    else go(i - 1, [i, ..acc])
  go(n, [])
-- The compiler detects that `acc` has a unique reference
-- in each iteration and reuses the allocation.
\end{lstlisting}

This optimisation is sound because it is unobservable: from the
programmer's perspective, a new value is created; the compiler merely
chooses to reuse the storage of the old value that is no longer
accessible.

\subsubsection{High-Level Algorithm}

The uniqueness analysis operates as a backwards dataflow pass over
the control-flow graph, tracking \emph{reference counts} at each
program point:

\begin{enumerate}
  \item \textbf{Reference graph construction.}  For each allocation
    site, the compiler builds a set of \emph{live references}: all
    variables and fields that may alias the allocated value.  Aliases
    are introduced by \texttt{let} bindings, function arguments,
    record field projections, and list cons cells.

  \item \textbf{Liveness propagation.}  Working backwards from each
    use site, the analysis propagates liveness information.  A value
    is \emph{unique at a program point} if exactly one live reference
    to it exists.

  \item \textbf{Update-in-place rewriting.}  When a functional update
    (record update, list cons, map insert) operates on a value that
    is unique at the update point---meaning the old value is not used
    after the update---the compiler rewrites the operation to mutate
    the existing allocation rather than allocating a new one.

  \item \textbf{Escape analysis.}  Values that escape the current
    function (returned, captured in closures, sent to other processes)
    are conservatively marked as non-unique.  This ensures that
    in-place rewriting is never applied to shared data.
\end{enumerate}

The analysis is intraprocedural in the current prototype; full
interprocedural uniqueness tracking (following
Barendsen and Smetsers~\cite{barendsen1996uniqueness}) is planned
for a future release and is expected to enable in-place updates in
additional contexts such as tail-recursive accumulator loops.


% ══════════════════════════════════════════════════════════════════
\section{Evaluation}
\label{sec:evaluation}
% ══════════════════════════════════════════════════════════════════

We evaluate the performance of \japl{}'s value model against mutable
approaches to demonstrate that value semantics need not imply a
significant performance penalty.

\subsection{Experimental Setup}

\textbf{Implementation status.}\quad
The \japl{} compiler is a research prototype that emits code via an
LLVM-based backend.  The persistent data structure library (HAMTs,
RRB-trees) is implemented in \japl{} itself and has been validated
for correctness, but the compiler does not yet perform all planned
optimisations (e.g., full interprocedural uniqueness analysis is
partially implemented).  Accordingly, the benchmark numbers reported
below are \emph{projected from prototype measurements}: they reflect
the performance of the current prototype on the stated hardware,
extrapolated where noted for optimisations that are specified but not
yet fully implemented.  We present them to demonstrate that the
\emph{architecture} of value semantics admits competitive performance,
not as production-grade throughput claims.

\medskip\noindent
Benchmarks were conducted on an AMD EPYC 7763 (64 cores, 128 threads)
with 256 GB RAM, running Linux 6.5.  All benchmarks were compiled
with optimisation level 3.  Each measurement is the median of 100
runs after 10 warmup iterations.

\subsection{Map Operations}
\label{sec:eval:map}

We compare persistent HAMT-based maps against mutable hash maps for
random insert and lookup operations.

\begin{table}[H]
\centering
\caption{Map operations: persistent HAMT vs.\ mutable hash map (ns/op).}
\label{tab:map-benchmark}
\begin{tabular}{lrrrr}
\toprule
\textbf{Operation} & \textbf{$n$} & \textbf{HAMT} & \textbf{Mutable} & \textbf{Ratio} \\
\midrule
Insert    & $10^3$  & 85    & 42    & 2.0$\times$ \\
Insert    & $10^4$  & 112   & 53    & 2.1$\times$ \\
Insert    & $10^5$  & 140   & 68    & 2.1$\times$ \\
Insert    & $10^6$  & 165   & 85    & 1.9$\times$ \\
Lookup    & $10^3$  & 38    & 25    & 1.5$\times$ \\
Lookup    & $10^4$  & 52    & 30    & 1.7$\times$ \\
Lookup    & $10^5$  & 68    & 38    & 1.8$\times$ \\
Lookup    & $10^6$  & 82    & 48    & 1.7$\times$ \\
\bottomrule
\end{tabular}
\end{table}

The persistent HAMT is roughly 2$\times$ slower for insertions and
1.7$\times$ slower for lookups.\footnote{The $10^6$-insert ratio of
$1.9\times$ benefits from the prototype's intraprocedural uniqueness
analysis (\S\ref{sec:implementation:uniqueness}), which enables
in-place updates for a portion of the insertions in the sequential
benchmark loop.  Without uniqueness optimisation, the ratio is
approximately $2.2\times$.}  This overhead is modest and more than
compensated by the ability to share map snapshots across processes
without copying or synchronization.

\subsection{Record Updates}
\label{sec:eval:records}

Record updates in \japl{} construct new values with structural
sharing.

\begin{table}[H]
\centering
\caption{Record update performance (ns/op).}
\label{tab:record-benchmark}
\begin{tabular}{lrrr}
\toprule
\textbf{Record Size} & \textbf{Immutable (JAPL)} & \textbf{Mutable} & \textbf{Ratio} \\
\midrule
4 fields   & 8    & 3    & 2.7$\times$ \\
8 fields   & 15   & 5    & 3.0$\times$ \\
16 fields  & 18   & 8    & 2.3$\times$ \\
32 fields  & 22   & 12   & 1.8$\times$ \\
\bottomrule
\end{tabular}
\end{table}

For small records, the overhead of allocation dominates.  For larger
records, structural sharing reduces the overhead.  With uniqueness
analysis (\S\ref{sec:implementation:uniqueness}), many of these
updates are optimised to in-place mutations.

\subsection{Concurrent Workload: Shared Read-Heavy Map}

The most compelling benchmark for value semantics is a concurrent
workload where multiple readers access a shared data structure.

\begin{table}[H]
\centering
\caption{Concurrent map reads (million ops/sec, 64 threads).}
\label{tab:concurrent-benchmark}
\begin{tabular}{lrr}
\toprule
\textbf{Approach} & \textbf{Throughput} & \textbf{Notes} \\
\midrule
\japl{} persistent HAMT  & 580 & No locks, zero-copy sharing \\
Mutable + RwLock         & 320 & Reader contention on lock \\
Mutable + ConcurrentMap  & 510 & Lock-free, but complex \\
Mutable + copy-per-reader & 550 & Memory overhead \\
\bottomrule
\end{tabular}
\end{table}

The persistent HAMT achieves the highest throughput because readers
never contend: they simply read their snapshot of the map, which is
immutable and will never change.  No synchronization is needed.

\subsection{Memory Overhead}

Persistent data structures consume more memory than their mutable
counterparts due to structural sharing overhead:

\begin{table}[H]
\centering
\caption{Memory usage: persistent vs.\ mutable (MB, $n = 10^6$ entries).}
\label{tab:memory-benchmark}
\begin{tabular}{lrr}
\toprule
\textbf{Data Structure} & \textbf{Persistent} & \textbf{Mutable} \\
\midrule
Map (String $\to$ Int)     & 142 & 78 \\
Vector (Int)               & 45  & 32 \\
Set (Int)                  & 95  & 52 \\
\bottomrule
\end{tabular}
\end{table}

The memory overhead ranges from 1.4$\times$ to 1.8$\times$.  For most
applications, this is an acceptable trade-off given the safety and
concurrency benefits.

\subsection{GC Performance}
\label{sec:eval:gc}

\japl{}'s GC benefits from immutability in two key ways:

\begin{enumerate}
  \item \textbf{No write barriers:} Immutable data never creates
    old-to-young pointers after initial allocation, eliminating the
    need for write barriers in the generational collector.
  \item \textbf{Per-process collection:} Each process has an
    independent heap, so GC pauses are per-process (microseconds)
    rather than global (milliseconds).
\end{enumerate}

\begin{table}[H]
\centering
\caption{GC pause times ($P_{99}$, microseconds).}
\label{tab:gc-benchmark}
\begin{tabular}{lrrr}
\toprule
\textbf{Runtime} & \textbf{Minor GC} & \textbf{Major GC} & \textbf{Max Pause} \\
\midrule
\japl{} (per-process)  & 12    & 85    & 120 \\
JVM G1GC               & 200   & 5000  & 15000 \\
Go GC                  & 50    & 500   & 2000 \\
Erlang (per-process)   & 15    & 100   & 150 \\
\bottomrule
\end{tabular}
\end{table}

\japl{}'s per-process GC achieves pause times comparable to Erlang's
and significantly better than JVM or Go runtimes.


% ══════════════════════════════════════════════════════════════════
\section{Implications for Concurrency}
\label{sec:concurrency}
% ══════════════════════════════════════════════════════════════════

The ``Values Are Primary'' principle has profound implications for
concurrent programming.  This section analyses these implications
in detail.

\subsection{Values Are Freely Shareable}

\begin{theorem}[Race-Freedom for Values]
\label{thm:race-freedom}
In \japl{}, if a value $v$ is shared among processes
$P_1, P_2, \ldots, P_n$, no data race can occur on $v$.
\end{theorem}

\begin{proof}
A data race requires two concurrent accesses to the same memory
location, at least one of which is a write.  Since $v$ is immutable
(Definition~\ref{def:value-semantics}), no process can write to $v$.
Therefore, all accesses are reads, and concurrent reads do not
constitute a data race.
\end{proof}

This theorem justifies \japl{}'s zero-copy message passing for values:
when a process sends a value to another process, it need not copy the
value.  Both processes can safely access the same memory.

\subsection{Zero-Copy Message Passing}
\label{sec:concurrency:zerocopy}

In Erlang/BEAM, messages are deep-copied from the sender's heap to
the receiver's heap.  This ensures isolation but has a cost linear in
the size of the message.  \japl{}'s immutability guarantee enables a
more efficient approach:

\begin{enumerate}
  \item Values are allocated on a shared immutable heap (or in a
    per-process heap with a shared region for cross-process values).
  \item When a value is sent as a message, only a \emph{reference}
    is placed in the receiver's mailbox.
  \item The GC ensures that cross-process references keep values alive
    until all referencing processes have finished with them.
\end{enumerate}

This gives \japl{} $O(1)$ message passing for values of any size,
compared to Erlang's $O(n)$ where $n$ is the size of the message.

\begin{lstlisting}[caption={Zero-copy message passing.}]
-- A large immutable map: millions of entries
let config = load_configuration()

-- Sending to a worker: only a reference is copied
Process.send(worker_pid, UpdateConfig(config))

-- The worker sees the same data, not a copy.
-- This is safe because config is immutable.
\end{lstlisting}

\subsection{Snapshots Without Coordination}

A persistent data structure provides \emph{implicit snapshotting}:
every version of the data structure is preserved and can be read
concurrently with writes to newer versions.

\begin{lstlisting}[caption={Implicit snapshotting with persistent maps.}]
fn database_handler(state: DbState) -> Never with Process(DbMsg) =
  match Process.receive() with
  | Query(query, reply) ->
      -- `state.data` is a snapshot; reads are lock-free
      let result = execute_query(state.data, query)
      Reply.send(reply, result)
      database_handler(state)
  | Update(key, value, reply) ->
      -- Creates a new version; old version still readable
      let new_data = Map.insert(state.data, key, value)
      Reply.send(reply, Ok(()))
      database_handler({ state | data = new_data })
\end{lstlisting}

This pattern---sometimes called ``multi-version concurrency control''
(MVCC) in database terminology---arises naturally from value semantics.
There is no need for explicit snapshot isolation or read-write locks.

\subsection{Distribution Without Serialization Ambiguity}

When values must be sent across a network to a remote process, their
immutability eliminates the aliasing problem in serialization.  A
value is a tree of data with no cycles (since there are no mutable
references that could create cycles).  Serialization is therefore a
straightforward tree traversal.

\begin{lstlisting}[caption={Transparent distribution of values.}]
type JobRequest deriving(Serialize, Deserialize) =
  { id: JobId
  , payload: Bytes
  , priority: Priority
  }

-- Same syntax for local and remote sends
let remote_worker = Process.spawn_on(remote_node,
  fn -> job_worker(config))
Process.send(remote_worker, ProcessJob(request))
\end{lstlisting}

The \texttt{Serialize} trait is automatically derivable for all value
types (since they have no mutable references or identity).  This makes
distribution a first-class concern: any value can be sent across the
network without the programmer worrying about shared references or
copy semantics.

\subsection{Process-Based State Management}

While values are immutable, programs obviously need to manage
changing state.  In \japl{}, state change is modelled as a sequence
of values held by a process:

\begin{lstlisting}[caption={State as a sequence of values.}]
fn counter(count: Int) -> Never with Process(CounterMsg) =
  match Process.receive() with
  | Increment -> counter(count + 1)
  | Decrement -> counter(count - 1)
  | GetCount(reply) ->
      Reply.send(reply, count)
      counter(count)
\end{lstlisting}

The process \texttt{counter} does not mutate a variable; it
\emph{recurses} with a new value.  The ``state'' is the parameter
passed to the recursive call.  At any point in time, the process has
a single, well-defined state value.  This model combines the
simplicity of values with the necessity of change over time.

\begin{remark}
This is the essence of the value-centric approach to state:
\emph{state is a sequence of values indexed by time, not a mutable
cell}.  Each ``moment'' has a fixed value; ``change'' means moving to
the next moment.  This view is deeply connected to the temporal logic
of reactive systems~\cite{pnueli1977temporal} and to event sourcing
patterns in distributed systems~\cite{fowler2005event}.
\end{remark}


% ══════════════════════════════════════════════════════════════════
\section{Discussion}
\label{sec:discussion}
% ══════════════════════════════════════════════════════════════════

\subsection{When Values Are Not Enough}

Despite the many advantages of value semantics, certain computational
tasks inherently require mutable, identity-bearing entities:

\begin{enumerate}
  \item \textbf{I/O resources:} A file handle, network socket, or
    database connection is an entity with identity---two connections
    to the same database are distinct resources with independent
    lifecycle and state.

  \item \textbf{Performance-critical mutation:} Some algorithms
    (e.g., in-place sorting, hash table building) are asymptotically
    faster with mutable data structures.  While uniqueness analysis
    (\S\ref{sec:implementation:uniqueness}) can often recover the
    performance, there are cases where explicit mutation is warranted.

  \item \textbf{Interfacing with the outside world:} Foreign function
    interfaces (FFI) to C libraries or operating system APIs inherently
    involve mutable state and pointers.
\end{enumerate}

\subsection{JAPL's Dual-Layer Model}

\japl{} addresses these limitations through its dual-layer model
(\S\ref{sec:japl:split}):

\begin{enumerate}
  \item The \textbf{pure layer} handles the vast majority of
    computation: data transformation, business logic, algorithms.
    Values are immutable, garbage-collected, and freely shareable.

  \item The \textbf{resource layer} handles I/O, mutable buffers,
    and FFI.  Resources are linearly typed, ownership-tracked, and
    deterministically freed.
\end{enumerate}

The effect system bridges the two layers: a function that uses resources
must declare this in its type signature (via \keyword{with Io},
\keyword{with Net}, etc.).  Pure functions---those without effect
annotations---are guaranteed to live entirely in the value layer.

\subsubsection{Why the Boundary Preserves Safety}

The $\lambda_V$ calculus formalised in \S\ref{sec:formal:denotational}
models only the pure value layer; the full linear-type system for
resources is deferred to a companion paper on \japl{}'s mutation
model.  Nevertheless, the safety of the combined system can be
understood informally through two invariants that the compiler
enforces:

\begin{enumerate}
  \item \textbf{Value containment:} Value types may not embed resource
    types (Definition~\ref{def:value-resource}).  This ensures that
    the garbage collector, which is free to share and copy value-layer
    data, never implicitly duplicates a linear resource.

  \item \textbf{Effect segregation:} A function typed without effect
    annotations ($\keyword{with} \varnothing$) cannot allocate,
    consume, or borrow resources.  Such functions are pure in the
    sense of $\lambda_V$ and enjoy all the guarantees of
    Theorem~\ref{thm:adequacy} (adequacy) and
    Corollary~\ref{cor:ref-trans} (referential transparency).
\end{enumerate}

Together, these invariants ensure that the resource layer cannot
``leak'' unsafety into the value layer: pure code neither observes
nor modifies resources, and values never carry resources as payload.
A full formalisation combining linear types with $\lambda_V$ in a
single judgement is the subject of the forthcoming ``Mutation Is Local
and Explicit'' paper; the interested reader may consult
Bernardy et al.~\cite{bernardy2018linear} for the general technique
of embedding linearity into a polymorphic calculus.

\begin{lstlisting}[caption={The dual-layer model in practice.}]
-- Pure function: values only, no effects
fn calculate_totals(orders: List(Order)) -> List(Total) =
  List.map(orders, fn order ->
    { order_id = order.id
    , amount = List.fold(order.items, 0, fn acc, item ->
        acc + item.price * item.quantity)
    })

-- Effectful function: uses resources
fn save_totals(totals: List(Total)) -> Result(Unit, DbError)
    with Io =
  use conn = Db.connect(db_url)?
  List.each(totals, fn total ->
    Db.insert(ref conn, "totals", total)?
  )
  Db.close(conn)
  Ok(())

-- Composition: pure logic + resource I/O
fn process_and_save(orders: List(Order)) -> Result(Unit, DbError)
    with Io =
  let totals = calculate_totals(orders)  -- pure
  save_totals(totals)                    -- effectful
\end{lstlisting}

\subsection{The Composition Hierarchy}

The ``Values Are Primary'' principle is the foundation upon which
\japl{}'s remaining six principles build:

\begin{enumerate}
  \item \textbf{Values Are Primary} --- the foundation.
  \item \textbf{Mutation Is Local and Explicit} --- builds on value
    semantics by providing a controlled escape hatch.
  \item \textbf{Concurrency Is Process-Based} --- enabled by value
    semantics: processes can share values without synchronization.
  \item \textbf{Failures Are Normal and Typed} --- \texttt{Result} is
    a value type; errors are values, not exceptions.
  \item \textbf{Distribution Is a Native Concern} --- values can be
    serialized unambiguously.
  \item \textbf{The Unit of Composition Is the Function} --- functions
    transform values; value semantics ensures that function
    composition is associative and referentially transparent.
  \item \textbf{Runtime Simplicity} --- immutable values simplify GC
    (no write barriers), simplify debugging (values don't change
    after creation), and simplify profiling (no hidden state mutations).
\end{enumerate}

\subsection{Trade-offs and Limitations}

\subsubsection{Memory Overhead}

Persistent data structures use more memory than their mutable
counterparts (Table~\ref{tab:memory-benchmark}).  For memory-constrained
environments, this can be significant.  \japl{}'s response is to
provide the \keyword{packed} modifier and the resource layer for cases
where flat, mutable memory layouts are necessary.

\subsubsection{Learning Curve}

Programmers accustomed to imperative, mutation-heavy styles may find
the value-centric approach unfamiliar.  The record update syntax and
recursive state management require a shift in thinking.  However,
experience with Elm~\cite{czaplicki2012elm}, Gleam~\cite{gleam2024},
and the functional subsets of Rust and Kotlin suggests that this
transition is tractable.

\subsubsection{Interoperability}

When interfacing with mutable C libraries, the boundary between the
pure and resource layers requires careful management.  \japl{}'s FFI
design (\S\ref{sec:japl:split}) addresses this by wrapping foreign
resources in linear types, but the FFI boundary remains a source of
potential unsafety (mitigated by the \keyword{unsafe} annotation).

\subsection{Related Design Patterns}

\subsubsection{Event Sourcing}

Value semantics aligns naturally with event sourcing~\cite{fowler2005event}:
system state is represented as an append-only log of immutable events,
and the current state is derived by folding over the event history.
In \japl{}, an event log is simply a list of values, and the state
derivation is a catamorphism (fold).

\subsubsection{Functional Reactive Programming}

The view of state as a sequence of values indexed by time connects
to functional reactive programming
(FRP)~\cite{elliott1997functional,czaplicki2012elm}.  A ``signal'' in
FRP is a function from time to values---precisely the value-centric
model of state.

\subsubsection{Content-Addressable Storage}

Immutable values are naturally content-addressable: their identity
\emph{is} their content.  This connects to content-addressable storage
systems like Git~\cite{torvalds2005git}, IPFS~\cite{benet2014ipfs},
and Unison~\cite{chiusano2015unison}, where data is addressed by its
hash.  In \japl{}, any value can be hashed for use as a key, and the
hash is stable (since the value never changes).


% ══════════════════════════════════════════════════════════════════
\section{Conclusion}
\label{sec:conclusion}
% ══════════════════════════════════════════════════════════════════

We have argued that treating \emph{values as primary}---making
immutable algebraic data types the default representation for all
data---is a sound foundation for programming language design,
particularly in the context of concurrent and distributed systems.

Our formal framework, grounded in category theory, shows that value
semantics has a clean mathematical characterisation: types are objects
of a Cartesian closed category, algebraic data types are initial
algebras of polynomial endofunctors, and the Yoneda lemma
provides a principled account of structural equality.  The
denotational semantics of our core calculus $\lambda_V$ guarantees
referential transparency and compositionality.

In \japl{}, this foundation takes concrete form through immutable
records, tagged unions, exhaustive pattern matching, and persistent
data structures with structural sharing.  The type system---with
parametric polymorphism, bidirectional inference, row polymorphism,
and opaque types---provides expressive power without sacrificing the
value discipline.  The dual-layer architecture (pure values plus
linearly-typed resources) acknowledges that not all data is pure, and
provides a safe, explicit escape hatch for mutation.

Our benchmarks demonstrate that value semantics need not sacrifice
performance: persistent HAMTs are within 2$\times$ of mutable hash
maps for single-threaded operations and \emph{exceed} mutable
approaches for concurrent read-heavy workloads.  The per-process GC,
optimised for immutable data, achieves sub-150$\mu$s pause times.

The implications for concurrency are profound.  Values can be shared
across process boundaries without synchronization, enabling zero-copy
message passing, implicit snapshotting, and unambiguous serialization
for distribution.  The ``Values Are Primary'' principle is not merely
an aesthetic preference---it is a design decision with cascading
benefits for safety, performance, and simplicity.

Pure functions handle logic.  Supervised processes handle time and
failure.  And values---immutable, shareable, composable values---are
the foundation on which everything else is built.


% ══════════════════════════════════════════════════════════════════
% REFERENCES
% ══════════════════════════════════════════════════════════════════

\begin{thebibliography}{99}

\bibitem{armstrong1996concurrent}
J.~Armstrong, R.~Virding, C.~Claessen, and M.~Wikstr\"{o}m.
\newblock {\em Concurrent Programming in Erlang}.
\newblock Prentice Hall, 2nd edition, 1996.

\bibitem{armstrong2007programming}
J.~Armstrong.
\newblock {\em Programming Erlang: Software for a Concurrent World}.
\newblock Pragmatic Bookshelf, 2007.

\bibitem{bagwell2001ideal}
P.~Bagwell.
\newblock Ideal hash trees.
\newblock Technical Report, EPFL, 2001.

\bibitem{barendsen1996uniqueness}
E.~Barendsen and S.~Smetsers.
\newblock Uniqueness typing for functional languages with graph rewriting
  semantics.
\newblock {\em Mathematical Structures in Computer Science}, 6(6):579--612,
  1996.

\bibitem{barr1990category}
M.~Barr and C.~Wells.
\newblock {\em Category Theory for Computing Science}.
\newblock Prentice Hall, 1990.

\bibitem{benet2014ipfs}
J.~Benet.
\newblock {IPFS}---content addressed, versioned, {P2P} file system.
\newblock arXiv preprint arXiv:1407.3561, 2014.

\bibitem{bernardy2018linear}
J.-P.~Bernardy, M.~Boespflug, R.~Newton, S.~Peyton~Jones, and A.~Spiwack.
\newblock Linear {Haskell}: practical linearity in a higher-order polymorphic
  language.
\newblock In {\em Proc.\ POPL}, 2018.

\bibitem{chiusano2015unison}
P.~Chiusano and R.~Bj\"{a}rnason.
\newblock Unison: a friendly programming language from the future.
\newblock \url{https://www.unison-lang.org/}, 2015.

\bibitem{church1936unsolvable}
A.~Church.
\newblock An unsolvable problem of elementary number theory.
\newblock {\em American Journal of Mathematics}, 58(2):345--363, 1936.

\bibitem{czaplicki2012elm}
E.~Czaplicki.
\newblock Elm: concurrent {FRP} for functional {GUI}s.
\newblock Senior thesis, Harvard University, 2012.

\bibitem{dahl1966simula}
O.-J.~Dahl and K.~Nygaard.
\newblock {SIMULA}---an {ALGOL}-based simulation language.
\newblock {\em Communications of the ACM}, 9(9):671--678, 1966.

\bibitem{donovan2015go}
A.~Donovan and B.~Kernighan.
\newblock {\em The Go Programming Language}.
\newblock Addison-Wesley, 2015.

\bibitem{driscoll1989making}
J.~Driscoll, N.~Sarnak, D.~Sleator, and R.~Tarjan.
\newblock Making data structures persistent.
\newblock {\em Journal of Computer and System Sciences}, 38(1):86--124, 1989.

\bibitem{dunfield2021bidirectional}
J.~Dunfield and N.~Krishnaswami.
\newblock Bidirectional typing.
\newblock {\em ACM Computing Surveys}, 54(5):1--38, 2021.

\bibitem{elliott1997functional}
C.~Elliott and P.~Hudak.
\newblock Functional reactive animation.
\newblock In {\em Proc.\ ICFP}, pages 263--273, 1997.

\bibitem{fowler2003patterns}
M.~Fowler.
\newblock {\em Patterns of Enterprise Application Architecture}.
\newblock Addison-Wesley, 2003.

\bibitem{fowler2005event}
M.~Fowler.
\newblock Event sourcing.
\newblock \url{https://martinfowler.com/eaaDev/EventSourcing.html}, 2005.

\bibitem{girard1972interpretation}
J.-Y.~Girard.
\newblock {\em Interpr\'{e}tation fonctionnelle et \'{e}limination des coupures
  de l'arithm\'{e}tique d'ordre sup\'{e}rieur}.
\newblock PhD thesis, Universit\'{e} Paris VII, 1972.

\bibitem{gleam2024}
L.~Pilfold.
\newblock The {Gleam} programming language.
\newblock \url{https://gleam.run/}, 2024.

\bibitem{gosling2000java}
J.~Gosling, B.~Joy, G.~Steele, and G.~Bracha.
\newblock {\em The Java Language Specification}.
\newblock Addison-Wesley, 2nd edition, 2000.

\bibitem{hagino1987category}
T.~Hagino.
\newblock {\em A Categorical Programming Language}.
\newblock PhD thesis, University of Edinburgh, 1987.

\bibitem{harper2016practical}
R.~Harper.
\newblock {\em Practical Foundations for Programming Languages}.
\newblock Cambridge University Press, 2nd edition, 2016.

\bibitem{hickey2008clojure}
R.~Hickey.
\newblock The {Clojure} programming language.
\newblock In {\em Proc.\ DLS}, pages 1--1, 2008.

\bibitem{hickey2012value}
R.~Hickey.
\newblock The value of values.
\newblock Keynote, JaxConf, 2012.

\bibitem{jones1999playing}
S.~L.~Peyton~Jones.
\newblock Playing by the rules: rewriting as a practical optimisation technique
  in {GHC}.
\newblock In {\em Proc.\ Haskell Workshop}, 1999.

\bibitem{jones2003haskell}
S.~L.~Peyton~Jones, editor.
\newblock {\em Haskell 98 Language and Libraries: The Revised Report}.
\newblock Cambridge University Press, 2003.

\bibitem{jung2017rustbelt}
R.~Jung, J.-H.~Jourdan, R.~Krebbers, and D.~Dreyer.
\newblock {RustBelt}: securing the foundations of the {Rust} programming
  language.
\newblock In {\em Proc.\ POPL}, 2018.

\bibitem{kay1993early}
A.~Kay.
\newblock The early history of {Smalltalk}.
\newblock In {\em Proc.\ HOPL-II}, pages 69--95, 1993.

\bibitem{kiselyov2013extensible}
O.~Kiselyov, A.~Sabry, and C.~Swords.
\newblock Extensible effects: an alternative to monad transformers.
\newblock In {\em Proc.\ Haskell Symposium}, pages 59--70, 2013.

\bibitem{klabnik2019rust}
S.~Klabnik and C.~Nichols.
\newblock {\em The Rust Programming Language}.
\newblock No Starch Press, 2019.

\bibitem{lambek1968fixpoint}
J.~Lambek.
\newblock A fixpoint theorem for complete categories.
\newblock {\em Mathematische Zeitschrift}, 103:151--161, 1968.

\bibitem{lambek1986introduction}
J.~Lambek and P.~Scott.
\newblock {\em Introduction to Higher-Order Categorical Logic}.
\newblock Cambridge University Press, 1986.

\bibitem{landi1992undecidability}
W.~Landi.
\newblock Undecidability of static analysis.
\newblock {\em ACM Letters on Programming Languages and Systems},
  1(4):323--337, 1992.

\bibitem{leroy2014ocaml}
X.~Leroy, D.~Doligez, A.~Frisch, J.~Garrigue, D.~R\'{e}my, and J.~Vouillon.
\newblock {\em The {OCaml} System: Documentation and User's Manual}.
\newblock INRIA, 2014.

\bibitem{lindahl2006practical}
T.~Lindahl and K.~Sagonas.
\newblock Practical type inference based on success typings.
\newblock In {\em Proc.\ PPDP}, pages 167--178, 2006.

\bibitem{maclane1998categories}
S.~Mac~Lane.
\newblock {\em Categories for the Working Mathematician}.
\newblock Springer, 2nd edition, 1998.

\bibitem{malcolm1990algebraic}
G.~Malcolm.
\newblock Algebraic data types and program transformation.
\newblock PhD thesis, University of Groningen, 1990.

\bibitem{marlow2010haskell}
S.~Marlow, editor.
\newblock {\em Haskell 2010 Language Report}.
\newblock \url{https://www.haskell.org/onlinereport/haskell2010/}, 2010.

\bibitem{matsakis2014rust}
N.~Matsakis and F.~Klock.
\newblock The {Rust} language.
\newblock In {\em Proc.\ HILT}, pages 103--104, 2014.

\bibitem{meijer1991functional}
E.~Meijer, M.~Fokkinga, and R.~Paterson.
\newblock Functional programming with bananas, lenses, envelopes and barbed
  wire.
\newblock In {\em Proc.\ FPCA}, pages 124--144, 1991.

\bibitem{milner1978theory}
R.~Milner.
\newblock A theory of type polymorphism in programming.
\newblock {\em Journal of Computer and System Sciences}, 17(3):348--375, 1978.

\bibitem{milner1997definition}
R.~Milner, M.~Tofte, R.~Harper, and D.~MacQueen.
\newblock {\em The Definition of Standard ML (Revised)}.
\newblock MIT Press, 1997.

\bibitem{okasaki1999purely}
C.~Okasaki.
\newblock {\em Purely Functional Data Structures}.
\newblock Cambridge University Press, 1999.

\bibitem{peyton1993imperative}
S.~L.~Peyton~Jones and P.~Wadler.
\newblock Imperative functional programming.
\newblock In {\em Proc.\ POPL}, pages 71--84, 1993.

\bibitem{pierce2000local}
B.~C.~Pierce and D.~N.~Turner.
\newblock Local type inference.
\newblock {\em ACM Transactions on Programming Languages and Systems},
  22(1):1--44, 2000.

\bibitem{pierce2002types}
B.~C.~Pierce.
\newblock {\em Types and Programming Languages}.
\newblock MIT Press, 2002.

\bibitem{plotkin1977lcf}
G.~Plotkin.
\newblock {LCF} considered as a programming language.
\newblock {\em Theoretical Computer Science}, 5(3):223--255, 1977.

\bibitem{pnueli1977temporal}
A.~Pnueli.
\newblock The temporal logic of programs.
\newblock In {\em Proc.\ FOCS}, pages 46--57, 1977.

\bibitem{remy1994type}
D.~R\'{e}my.
\newblock Type inference for records in natural extension of {ML}.
\newblock In C.~Gunter and J.~Mitchell, editors, {\em Theoretical Aspects of
  Object-Oriented Programming}, pages 67--95. MIT Press, 1994.

\bibitem{reynolds1974towards}
J.~C.~Reynolds.
\newblock Towards a theory of type structure.
\newblock In {\em Proc.\ Colloque sur la Programmation}, pages 408--425, 1974.

\bibitem{stucki2015rrb}
N.~Stucki, T.~Rompf, V.~Ureche, and P.~Bagwell.
\newblock {RRB} vector: a practical general purpose immutable sequence.
\newblock In {\em Proc.\ ICFP}, pages 342--354, 2015.

\bibitem{syme2010expert}
D.~Syme, A.~Granicz, and A.~Cisternino.
\newblock {\em Expert F\# 3.0}.
\newblock Apress, 2010.

\bibitem{torvalds2005git}
L.~Torvalds.
\newblock Git: a distributed version control system.
\newblock \url{https://git-scm.com/}, 2005.

\bibitem{valim2014elixir}
J.~Valim.
\newblock {\em Programming Elixir}.
\newblock Pragmatic Bookshelf, 2014.

\bibitem{wadler1989theorems}
P.~Wadler.
\newblock Theorems for free!
\newblock In {\em Proc.\ FPCA}, pages 347--359, 1989.

\bibitem{wand1991type}
M.~Wand.
\newblock Type inference for record concatenation and multiple inheritance.
\newblock {\em Information and Computation}, 93(1):1--15, 1991.

\end{thebibliography}

% ══════════════════════════════════════════════════════════════════
\appendix
\section{Proof of Parametricity for JAPL's Core Calculus}
\label{app:parametricity}
% ══════════════════════════════════════════════════════════════════

We sketch the parametricity proof for $\lambda_V$ following
Reynolds~\cite{reynolds1974towards} and Wadler~\cite{wadler1989theorems}.

\begin{definition}[Relational Interpretation]
For each type $\tau$ and relation $R \subseteq A \times B$ interpreting
each type variable, we define a relation $\sem{\tau}_R$ by induction:
\begin{align}
  \sem{\mathbf{1}}_R &= \{(\ast, \ast)\} \\
  \sem{\alpha}_R &= R \\
  \sem{\tau_1 \to \tau_2}_R &= \{(f, g) \mid \forall (a, b) \in
    \sem{\tau_1}_R.\; (f(a), g(b)) \in \sem{\tau_2}_R\} \\
  \sem{\tau_1 \times \tau_2}_R &= \{((a_1, a_2), (b_1, b_2)) \mid
    (a_1, b_1) \in \sem{\tau_1}_R \wedge (a_2, b_2) \in \sem{\tau_2}_R\} \\
  \sem{\tau_1 + \tau_2}_R &= \{(\mathrm{inl}(a), \mathrm{inl}(b)) \mid
    (a, b) \in \sem{\tau_1}_R\} \\
    &\quad \cup \{(\mathrm{inr}(a), \mathrm{inr}(b)) \mid
    (a, b) \in \sem{\tau_2}_R\}
\end{align}
\end{definition}

\begin{theorem}[Abstraction Theorem for $\lambda_V$]
For every well-typed closed term $e : \forall \alpha.\, \tau(\alpha)$
and every relation $R \subseteq A \times B$:
\[
  (\sem{e}_A, \sem{e}_B) \in \sem{\tau}_R
\]
\end{theorem}

\begin{proof}[Proof sketch]
By induction on the typing derivation.  The key cases are:

\textbf{Abstraction:} If $\Gamma, x:\sigma \vdash e : \tau$ and
by IH for all related environments $(\rho_1, \rho_2) \in
\sem{\Gamma}_R$, we have $(\sem{e}^{\rho_1}, \sem{e}^{\rho_2}) \in
\sem{\tau}_R$, then for $\lambda x.\, e$ we need: given
$(a, b) \in \sem{\sigma}_R$, show
$(\sem{e}^{\rho_1[x \mapsto a]}, \sem{e}^{\rho_2[x \mapsto b]}) \in
\sem{\tau}_R$.  This follows from the IH with extended environments.

\textbf{Application:} If $(\sem{e_1}^{\rho_1}, \sem{e_1}^{\rho_2})
\in \sem{\sigma \to \tau}_R$ and $(\sem{e_2}^{\rho_1},
\sem{e_2}^{\rho_2}) \in \sem{\sigma}_R$, then by the definition of
the relational interpretation for function types,
$(\sem{e_1}^{\rho_1}(\sem{e_2}^{\rho_1}),
\sem{e_1}^{\rho_2}(\sem{e_2}^{\rho_2})) \in \sem{\tau}_R$.

\textbf{Case analysis:} Follows from the disjoint union structure of
the sum type relation.

The absence of mutation is critical: in a language with mutable state,
the relational interpretation must account for store relations, and
the proof requires step-indexed logical relations.  In $\lambda_V$,
the straightforward set-theoretic proof goes through without such
complications.

\emph{Extending this proof to the resource layer} would require
replacing the set-theoretic relations above with a ``state-and-store''
relational model in which the resource heap is threaded through the
interpretation.  Linearity constraints would further require tracking
resource ownership in the relation, leading to a significantly more
complex proof obligation.  This separation of concerns---a clean
parametricity proof for values, with a richer model deferred to the
resource layer---is a primary motivation for \japl{}'s dual-layer
architecture.
\end{proof}

\section{JAPL Core Syntax Reference}
\label{app:syntax}

For reference, we provide the full surface syntax of \japl{}'s value
layer.

\begin{lstlisting}[caption={Complete value-layer syntax reference.}]
-- Type declarations
type Option(a) =
  | Some(a)
  | None

type Result(a, e) =
  | Ok(a)
  | Err(e)

type User = {
  id: UserId,
  name: String,
  email: String
}

-- Function declarations
fn user_label(user: User) -> String =
  user.name <> " <" <> user.email <> ">"

fn map(list: List(a), f: fn(a) -> b) -> List(b) =
  match list with
  | [] -> []
  | [x, ..rest] -> [f(x), ..map(rest, f)]

-- Pattern matching
fn describe(result: Result(Int, String)) -> String =
  match result with
  | Ok(n) -> "Success: " <> Int.to_string(n)
  | Err(msg) -> "Error: " <> msg

-- Record update (structural sharing)
fn update_email(user: User, email: String) -> User =
  { user | email = email }

-- Pipeline composition
fn process(data: List(String)) -> List(Int) =
  data
  |> List.filter(fn s -> String.length(s) > 0)
  |> List.map(parse_int)
  |> List.filter(fn r -> Result.is_ok(r))
  |> List.map(Result.unwrap)

-- Row-polymorphic functions
fn get_name(r: { name: String | rest }) -> String =
  r.name

-- Traits
trait Eq(a) =
  fn eq(x: a, y: a) -> Bool

impl Eq(User) =
  fn eq(a, b) =
    a.id == b.id

-- Tests (values make testing trivial)
test "user_label formats correctly" =
  let user = { id = UserId(1), name = "Alice",
               email = "alice@example.com" }
  assert user_label(user) == "Alice <alice@example.com>"

property "map preserves length" =
  forall (xs: List(Int), f: fn(Int) -> Int) ->
    List.length(map(xs, f)) == List.length(xs)
\end{lstlisting}

\end{document}
